% Requires amsmath, amssymb, and amsthm.
\section{Channel-specific delays for finite Gaussian records}\label{sec:universal}
\newcommand{\MC}{\mathcal C}
\newcommand{\MD}{\mathcal D}
\newcommand{\MT}{\mathbb T}
\newcommand{\MR}{\mathbb R}
\newcommand{\MS}{\mathbb S}

A finite Gaussian law is determined by its retained means and covariances. The construction below matches those quantities exactly with a reversible source and channel-specific delays. The source and its unobserved extension can depend on the target law.

The difficulty is to satisfy reversibility and exact covariance matching at the same time. A real source spectrum has no complex channel phases of its own. Separated delays can approximate the phases of each rank-one spectral mass, but an approximation does not give the same finite law. We therefore apply the approximation to nearby valid covariance arrays that surround the target. Small changes preserve this containment, so a positive mixture of their outputs matches the target exactly. This also permits smoothing the masses into a strictly positive density before taking the final mixture. The notation below specifies the finite set of moments that must be preserved.

Let $m,N\ge1$ be integers. A retained covariance array consists of real
numbers $C_{ij}(k)$, $1\le i,j\le m$, $|k|\le N-1$, with
$C_{ji}(-k)=C_{ij}(k)$. Its block Toeplitz matrix is
$\Gamma_N(C)=(C_{ij}(t-u))_{(t,i),(u,j)}$, $0\le t,u<N$.
The cone $\MC_{m,N}$ contains the arrays for which $\Gamma_N(C)\succeq0$.
Use the independent coordinates $C_{ii}(k)$, $0\le k<N$, and
$C_{ij}(k)$, $i<j$, $|k|<N$. Their real dimension is
\begin{equation}\label{eq:mc-dimension}
 d_{\mathrm{mom}}=mN+\binom m2(2N-1)=m^2N-\frac{m(m-1)}2.
\end{equation}
Throughout, $\|C\|_{\mathrm{coord}}$ is the maximum absolute value of
these independent coordinates. Each entry of $\Gamma_N(C)$ is one
coordinate, possibly repeated. In particular
\begin{equation}\label{eq:mc-entry-norm}
 \|\Gamma_N(C)\|_{\mathrm{op}}\le mN\|C\|_{\mathrm{coord}}.
\end{equation}
Spectral densities are with respect to $\nu(d\omega)=d\omega/(2\pi)$
on $\MT=\MR/(2\pi\mathbb Z)$, with moment convention
$C_{ij}(k)=\int e^{ik\omega}\mu_{ij}(d\omega)$.
For integer channel delays $d=(d_1,\ldots,d_m)$ put
\[
 D_d(\omega)=\operatorname{diag}(e^{-id_1\omega},\ldots,e^{-id_m\omega}).
\]
Thus $Y_i(t)=X_i(t-d_i)$ has spectral measure $D_d\mu D_d^*$ and
\begin{equation}\label{eq:mc-delay-sign}
 C^Y_{ij}(k)=C^X_{ij}(k-d_i+d_j).
\end{equation}

\begin{theorem}[Common hierarchical delays]\label{thm:mc-main}
Let $C^{(1)},\ldots,C^{(s)}\in\MC_{m,N}$ be a nonempty finite collection
with $\Gamma_N(C^{(c)})\succ0$. Write
$\lambda_c=\lambda_{\min}(\Gamma_N(C^{(c)}))$ and
$v_c=\max_i C^{(c)}_{ii}(0)$.
If $m\ge2$, every integer $M\ge2$ satisfying
\begin{equation}\label{eq:mc-threshold}
 M-1>\pi(N+1)(2d_{\mathrm{mom}}-1)
       \max_{1\le c\le s}\left(\frac{2mN v_c}{\lambda_c}+1\right)
\end{equation}
admits the following simultaneous realizations with the shared delay vector
\begin{equation}\label{eq:mc-delays}
 (d_1,\ldots,d_m)=(M,M^2,\ldots,M^{m-1},0).
\end{equation}
For each $c$ there is a real stationary reversible Gaussian process
$X^{(c)}$, whose spectral density $S_c$ is a real, even, symmetric matrix
trigonometric polynomial, such that $S_c(\omega)\succeq\epsilon I_m$ for
all $\omega$ and a common $\epsilon>0$. The delayed process
$Y^{(c)}_i(t)=X^{(c)}_i(t-d_i)$ has retained covariance exactly $C^{(c)}$.
The density degrees have a common finite upper bound $L$, and each
source is $L$-dependent. For every coordinate,
\begin{equation}\label{eq:mc-variance}
 \operatorname{Var}(X^{(c)}_i(t))=C^{(c)}_{ii}(0).
\end{equation}
Consequently these are exact realizations of the retained Gaussian laws,
with zero added errors. Arbitrary constant mean vectors are preserved by
assigning the same mean vectors to the sources. Each delay response has
unit coefficient norm.

For $m=1$ the same conclusion holds with $d_1=0$, without a restriction
on $M$. For $N=1$ it holds for every prescribed vector of nonnegative
integer delays, without condition~\eqref{eq:mc-threshold}.
\end{theorem}

\begin{proof}
We first reduce the retained law to finitely many rank-one spectral masses. We then move their frequencies so the prescribed delays approximate their phases. Finally, convex mixtures correct the retained moments without a signed spectral perturbation that could destroy positivity.

\emph{Finite spectral representation.}
Every $C\in\MC_{m,N}$ has a finite positive semidefinite matrix spectral
measure obeying $\mu(-A)=\overline{\mu(A)}$ and the prescribed retained
moments. To see this, take a real Gram representation
$C_{ij}(t-u)=\langle v_{t,i},v_{u,j}\rangle$ in a real space $E$ of
dimension $\operatorname{rank}\Gamma_N(C)$. The map
$v_{t,i}\mapsto v_{t+1,i}$ on the span with $0\le t<N-1$ is well defined
and isometric: Toeplitz equality preserves the squared norm of every
linear combination, including a zero combination. The domain and range have equal dimension. Thus extend this map to a
real orthogonal map $U:E\to E$. When $N=1$, one may simply take $U=I_E$. Then
$v_{t,i}=U^t v_{0,i}$ for the retained times. The spectral decomposition
of the complexification of $U$ supplies the measure. More explicitly,
with a complex inner product linear in its first argument, let
$P_\omega$ be the eigenspace projection for $e^{i\omega}$ and set
$W_{\omega,ij}=\langle P_\omega v_{0,i},P_\omega v_{0,j}\rangle$.
These are Hermitian positive semidefinite matrices and
$C_{ij}(k)=\sum_\omega e^{ik\omega}W_{\omega,ij}$.
Realness of $U$ pairs conjugate eigenspaces. The endpoint matrices at
$0$ and $\pi$ are real. There are at most $\dim E\le mN$ distinct nodes.

Equivalently, write the real moments as a finite sum
\begin{equation}\label{eq:mc-paired-atoms}
 C_{ij}(k)=\sum_\ell\operatorname{Re}
                (e^{ik\omega_\ell}v_{\ell,i}\overline{v_{\ell,j}}).
\end{equation}
Here a conjugate pair of matrix masses has first been combined into its
half-circle mass and then split into rank-one matrices; endpoint masses
can be split into real rank-one matrices. Each term in
\eqref{eq:mc-paired-atoms} represents the full-circle measure
$(vv^*\delta_\omega+\overline v v^{\mathsf T}\delta_{-\omega})/2$.
In particular,
\begin{equation}\label{eq:mc-rankone-budget}
 \sum_\ell |v_{\ell,i}|^2=C_{ii}(0),\qquad
 \sum_\ell |v_{\ell,i}||v_{\ell,j}|
 \le\sqrt{C_{ii}(0)C_{jj}(0)}.
\end{equation}

\emph{Simultaneous phase approximation.}
The retained moments change little under a small frequency displacement. The hierarchical delays make such a displacement sufficient to set one channel phase at a time, with decreasing effects on phases already set.
Fix $m\ge2$, $M\ge2$, an arbitrary real representative $\omega$ of a
frequency, and phases $p_1,\ldots,p_{m-1}$ on the unit circle. Starting
with $\theta_0=\omega$, choose $\theta_j$ nearest to $\theta_{j-1}$ among
the real solutions of $e^{-iM^j\theta_j}=p_j$. Their spacing is
$2\pi/M^j$, so $|\theta_j-\theta_{j-1}|\le\pi/M^j$. Put
$\theta=\theta_{m-1}$. Then
\begin{align}
 |\theta-\omega|&\le\sum_{j=1}^{m-1}\frac\pi{M^j}
                         \le\frac\pi{M-1},\label{eq:mc-grid-freq}\\
 |e^{-iM^i\theta}-p_i|
 &\le M^i|\theta-\theta_i|
 \le\pi\sum_{r=1}^{m-1-i}M^{-r}
 \le\frac\pi{M-1},\quad 1\le i<m.\label{eq:mc-grid-phase}
\end{align}
The sum in the last line is empty for $i=m-1$, whose phase is exact.
Frequencies are reduced modulo $2\pi$ only after this construction;
no restriction to a half-circle is imposed on $\theta$.

Apply this construction separately to each vector $v$ in
\eqref{eq:mc-paired-atoms}. Multiplying $v$ by a global unit complex
number does not change $vv^*$. If $v_m\ne0$, use this freedom to make
$v_m$ real nonnegative. If $v_m=0$, use any global phase. Set
$a_i=|v_i|$ and write $v_i=p_i a_i$, taking $p_m=1$ and assigning any
unit phase to a zero coordinate. Thus $v=\operatorname{diag}(p)a$ with
$a\in[0,\infty)^m$ even when some coordinates vanish. Give the source
the real-even measure
\[
 \frac12 aa^{\mathsf T}\delta_\theta+
 \frac12 aa^{\mathsf T}\delta_{-\theta}.
\]
Its delayed moment contribution is
$\operatorname{Re}(e^{ik\theta}q_i\overline q_j)a_i a_j$, where
$q_i=e^{-id_i\theta}$ and $q_m=1$. At a coincident reflected pair
($\theta=0$ or $\pi$ modulo $2\pi$), the two half masses add to one
mass, and every $q_i$ is real because $d_i$ is integral. The output mass
there is therefore real and admissible. At every other frequency the
reflected output mass is the conjugate of the first one. These
observations also cover a frequency displaced across an endpoint.

For $|k|\le N-1$, equations~\eqref{eq:mc-grid-freq}--\eqref{eq:mc-grid-phase}
give
\begin{align*}
 |e^{ik\theta}q_i\overline q_j-e^{ik\omega}p_i\overline p_j|
 &\le |k||\theta-\omega|+|q_i-p_i|+|q_j-p_j|\\
 &\le\frac{\pi(N+1)}{M-1}.
\end{align*}
Combining all the real source atoms gives an output $\widetilde C$ with
\begin{equation}\label{eq:mc-approx}
 \|\widetilde C-C\|_{\mathrm{coord}}
 \le\frac{\pi(N+1)V}{M-1},\qquad
 \widetilde C_{ii}(0)=C_{ii}(0),\quad
 V=\max_i C_{ii}(0),
\end{equation}
by~\eqref{eq:mc-rankone-budget}. The same delay vector works for every
such array. Splitting into rank-one terms before moving frequencies is
essential; no simultaneous diagonal dephasing of a general matrix mass
has been assumed.

\emph{A simplex with an explicit perturbation margin.}
The approximate output need not equal the target. To correct this error while retaining positive source measures, we approximate several arrays around the target. A simplex is their convex hull. The margin below ensures that the target stays inside that hull after the approximation.
For a target $C^{(c)}$, use the natural coordinate basis $E_1,\ldots,E_{d_{\mathrm{mom}}}$
and write
\[
 a_c=\frac{\lambda_c}{2mN},\quad
 B_c=v_c+a_c,\quad
 V_{c,0}=C^{(c)}-a_c\sum_{j=1}^{d_{\mathrm{mom}}} E_j,\quad
 V_{c,j}=C^{(c)}+a_cE_j\quad(1\le j\le d_{\mathrm{mom}}).
\]
Every matrix perturbation has operator norm at most $mNa_c=\lambda_c/2$
by~\eqref{eq:mc-entry-norm}; hence all vertices are positive definite.
Their maximal coordinate variances are at most $B_c$.
This simplex contains the coordinate cube with center $C^{(c)}$ and
radius
\begin{equation}\label{eq:mc-radius}
 \rho_c=\frac{a_c}{2d_{\mathrm{mom}}-1}.
\end{equation}
Indeed the barycentric coordinates of displacement $x$ are
\[
 \alpha_0=\frac{1-\sum_j x_j/a_c}{d_{\mathrm{mom}}+1},\qquad
 \alpha_j=\alpha_0+x_j/a_c.
\]
For $\|x\|_\infty\le\rho_c$, the numerator of $\alpha_0$ is at least
$1-d_{\mathrm{mom}}/(2d_{\mathrm{mom}}-1)\ge0$, and the numerator of $\alpha_j$ is at least
$1-(2d_{\mathrm{mom}}-1)/(2d_{\mathrm{mom}}-1)=0$. These inequalities also hold for $d_{\mathrm{mom}}=1$.

If each vertex of this simplex moves by at most $\eta<\rho_c$ in
coordinate norm, its new convex hull contains the cube of radius
$\rho_c-\eta$ about $C^{(c)}$. For each linear functional $u$, the hull's support function drops by at most
$\eta\|u\|_1$. The original cube provides support
$u\cdot C^{(c)}+\rho_c\|u\|_1$. The
support-function characterization of containment proves the claim.
In particular the perturbed hull is full dimensional. Since it has
$d_{\mathrm{mom}}+1$ vertices, they are affinely independent and $C^{(c)}$ has strictly
positive barycentric coordinates.

Apply~\eqref{eq:mc-approx} to the vertices. The coordinate errors are at
most $\eta_c=\pi(N+1)B_c/(M-1)$. Condition~\eqref{eq:mc-threshold} is
exactly $\eta_c<\rho_c$ for each $c$. Thus every target belongs to the
strict interior of a simplex of outputs of finite reversible source
measures under the prescribed delay vector. The strict inequality leaves a positive margin for the next step. Without that margin, smoothing could move the output vertices so that their convex hull no longer contains the target.

\emph{Positive polynomial source densities at fixed delays.}
The convex hull currently contains outputs of atomic source measures. Smoothing gives polynomial densities and adding white power makes them strictly positive. Both operations change the moments, so they must fit inside the remaining containment margin before the exact mixture is chosen.
Now hold $M$ and the delay vector fixed. Denote the source measure at a
perturbed vertex by $\mu_{c,j}$. For $L\ge0$ let
\[
 F_L(\omega)=\frac1{L+1}\left|\sum_{r=0}^L e^{ir\omega}\right|^2
   =\sum_{|r|\le L}\left(1-\frac{|r|}{L+1}\right)e^{ir\omega}.
\]
This is a real, even, nonnegative trigonometric polynomial of integral
one against $\nu$. For $\epsilon>0$ define
\begin{equation}\label{eq:mc-smooth}
 S_{c,j}^{L,\epsilon}(\omega)
   =\int F_L(\omega-\theta)\mu_{c,j}(d\theta)+\epsilon I_m.
\end{equation}
These are real-even symmetric polynomial densities bounded below by
$\epsilon I_m$. For each fixed integer $r$, convolution multiplies the
$r$th source moment by $1-|r|/(L+1)$ when $L\ge|r|$. Added white power
changes only the zero-lag diagonal moments. By~\eqref{eq:mc-delay-sign},
only source lags of absolute value at most
\begin{equation}\label{eq:mc-Q}
 Q=N-1+M^{m-1}
\end{equation}
enter the retained output moments. Thus the smoothed output vertices
converge to their unsmoothed counterparts as $L\to\infty$ and
$\epsilon\to0$. There are finitely many vertices, so finite common
$L$ and positive common $\epsilon$ preserve all strict simplex
containments. With their resulting positive barycentric coefficients
$\beta_{c,j}$, put $S_c=\sum_{j=0}^{d_{\mathrm{mom}}}\beta_{c,j}S_{c,j}^{L,\epsilon}$.
The coefficients sum to one, so $S_c\succeq\epsilon I_m$, and linearity
of the output moment map gives exactly $C^{(c)}$.

These matrix densities define stationary real Gaussian processes.
Their source covariance matrices are real symmetric and even in the
lag, which makes their Gaussian finite-dimensional distributions
invariant under time reversal. Polynomial degree at most $L$ makes
all source covariances at $|k|>L$ vanish. Joint Gaussianity then gives
independence for any finite sets of source variables separated in time
by more than $L$; a monotone-class argument gives independence of the
corresponding generated sigma fields. Thus the sources are $L$-dependent.
At every frequency, $D_d$ is diagonal unitary and hence leaves the
diagonal spectral measures unchanged. Exact matching of the retained
zero-lag diagonal moments therefore proves~\eqref{eq:mc-variance}.

\emph{Scalar records and records of length one.}
For $m=1$, the finite spectral representation has an even scalar
measure and already gives a reversible source at every
simplex vertex without a phase approximation or delay. The same
smoothing and simplex correction complete the argument.
For $N=1$ and any fixed prescribed integer delay vector, represent each
positive definite simplex vertex by its matrix mass at frequency zero.
All delay factors are one there, so the vertex is already represented
exactly by a real-even source measure. Holding the delays fixed,
smoothing and simplex correction again give the claimed conclusion.
For zero delays alone, white sources suffice directly.
\end{proof}

\begin{proposition}[Uniform bounds at a fixed chosen delay]\label{prop:mc-uniform}
Fix $m\ge2$, $N\ge1$, $\lambda_*>0$ and $v_*<\infty$. The preceding
conclusion applies to every target satisfying
$\Gamma_N(C)\succeq\lambda_*I_{mN}$ and $\max_iC_{ii}(0)\le v_*$,
using a common delay parameter $M$ whenever
\[
 M\ge2,\qquad M-1>\pi(N+1)(2d_{\mathrm{mom}}-1)
                          \left(\frac{2mN v_*}{\lambda_*}+1\right).
\]
The targets may be an infinite family; their sources may differ. In
particular this applies uniformly to every compact subset of the
positive definite moment cone. For each fixed such $M$, the sources
can also have common finite polynomial degree and a common positive
spectral lower bound. One sufficient explicit choice is as follows:
\begin{gather*}
 a=\frac{\lambda_*}{2mN},\quad B=v_*+a,\quad
 \rho=\frac a{2d_{\mathrm{mom}}-1},\quad
 \eta=\frac{\pi(N+1)B}{M-1},\quad b=\rho-\eta>0,\\
 Q=N-1+M^{m-1},\qquad \epsilon=\frac b3,\qquad
 L\ge Q,\qquad L+1>\frac{3QB}{b}.
\end{gather*}
Every source obtained by the construction then obeys
$S_C(\omega)\succeq\epsilon I_m$ and has degree at most $L$.
These are sufficient bounds, not minimum resource requirements; they
depend on the fixed chosen $M$.
\end{proposition}
\begin{proof}
Use the same $a,B,\rho$ for the simplex about every target. The atomic
approximation preserves each source diagonal mass, so the source
measure at every vertex has coordinate variances at most $B$. For any
integer $r$, its matrix moment satisfies
\[
 |\widehat\mu_{ij}(r)|\le
   \sqrt{\widehat\mu_{ii}(0)\widehat\mu_{jj}(0)}\le B.
\]
For the finite real source measures constructed above, this follows
directly by summing the absolute values of their rank-one contributions
and applying Cauchy--Schwarz. At a required source lag $|r|\le Q$,
Fej\'er smoothing therefore changes a moment by at most $QB/(L+1)$.
The added white density changes each output coordinate by at most
$\epsilon$; for diagonal delays it affects only a zero-lag diagonal
coordinate. The total smoothed vertex displacement from the original
simplex is consequently at most
\[
 \eta+\frac{QB}{L+1}+\epsilon
 <\eta+\frac b3+\frac b3<\rho.
\]
The simplex argument applies separately to every target with the same
$M,L,\epsilon$. Compactness is not otherwise needed: a compact subset
of the positive definite cone supplies the displayed uniform eigenvalue
and variance bounds by continuity.
\end{proof}

\paragraph{Scope.}
The theorem concerns exact finite Gaussian record laws. It permits the
source law to vary with the target and with the selected delays. It does
not assert equality to a prescribed full-process law, a smallest delay,
a source range independent of $M$, or equality of non-Gaussian laws
from covariance matching. Its maximal constructed delay $M^{m-1}$ is
a sufficient choice. Polynomial source spectra make every covariance vanish beyond a finite lag.

\subsection{Limits of the finite-record construction}
The rank-one decomposition cannot in general be replaced by dephasing
a whole matrix mass. For example, with $r=2/5$ the Hermitian matrix
\[
 W=\begin{pmatrix}1&r&-ir\\r&1&r\\ir&r&1\end{pmatrix}
\]
has least eigenvalue at least $1-2r=1/5$. Its cycle product
$W_{12}W_{23}W_{31}=ir^3$ is invariant under diagonal unitary
congruence. A real symmetric matrix has a real cycle product, so no
such congruence makes $W$ real. Moving its rank-one pieces separately
avoids this obstruction.

Strict positive definiteness of the retained target is necessary for
the asserted positive source floor. If $S_X\succeq\epsilon I_m$, then
$D_dS_XD_d^*\succeq\epsilon I_m$. For any block vector $u$, integration
of its Fourier polynomial and Parseval's identity give
$u^{\mathsf T}\Gamma_N(C^Y)u\ge\epsilon\|u\|_2^2$.
A singular retained covariance therefore cannot have this type of
realization. This does not exclude realizations with singular spectra.

A prescribed full-process spectrum imposes further restrictions.
Consider a bivariate spectrum with unit diagonal and cross entry
\[
 f(\omega)=\frac14\left(1+\frac i2\sin\omega\right).
\]
It has the real-process conjugation symmetry and is uniformly positive,
since $|f(\omega)|\le\sqrt5/8<1$. A reversible-source explanation by
integer delays would require $e^{i\Delta\omega}f(\omega)$ to be real,
where $\Delta=d_1-d_2\in\mathbb Z$. The derivative of its imaginary
part at zero is $(\Delta+1/2)/4$, which never vanishes for integer
$\Delta$. Hence this full-process law has no such explanation, although
every finite retained covariance is positive definite and satisfies
Theorem~\ref{thm:mc-main}. The source and the unobserved output extension
may depend on the chosen finite record.

The source-degree bound in Proposition~\ref{prop:mc-uniform} also must
depend on the selected delay. For $m\ge2$, fix a temporally white
positive definite target whose contemporaneous covariance is the identity except for
$C_{1m}(0)=C_{m1}(0)=r$, with $0<|r|<1$. Under the hierarchical delays,
exactness requires $C^X_{1m}(-M)=r$. A polynomial source spectrum of
degree $L$ has zero moments at all lags of magnitude greater than $L$,
so necessarily $L\ge M$. Thus no bound independent of $M$ is possible,
even for a singleton compact family of targets. The bounds uniform over
targets in Proposition~\ref{prop:mc-uniform} hold at each fixed $M$.

\subsection{A linear-delay construction under an observable covariance margin}\label{sec:linear-delay}
A stronger condition on the retained covariances permits an explicit construction with linearly spaced delays. The comparison matrix below combines lower bounds for the auto spectra with absolute bounds for the cross spectra. Its positive definiteness is sufficient to keep the constructed source spectrum positive. This gives a direct construction for a smaller covariance class than the universal theorem.

For each condition $c$, put $H=N-1$, $v_{ci}=C^{(c)}_{ii}(0)$ and define
\[
a_{ci}=v_{ci}-2\sum_{k=1}^{H}|C^{(c)}_{ii}(k)|,\qquad
(B_c)_{ij}=2\sum_{k=-H}^{H}|C^{(c)}_{ij}(k)|\quad(i\ne j),
\]
with $(B_c)_{ii}=0$. Let
\[
Q_c=\operatorname{diag}(a_{c1},\ldots,a_{cm})-B_c.
\]
Assume $N,m\ge2$ and $Q_c\succ0$ in every condition. This sufficient condition uses only retained covariances. Positive definiteness of the observed record alone does not imply this condition. The simpler inequalities $a_{ci}>\sum_{j\ne i}(B_c)_{ij}$ imply $Q_c\succ0$ by strict diagonal dominance.

\begin{theorem}[Direct positive spectra with linearly spaced delays]\label{thm:linear-delay}
For any nonempty finite family satisfying the displayed condition, the common
channel delays
\[
 d_i=(m-i)(N-1),\qquad 1\le i\le m,
\]
realize all retained laws exactly with zero errors and reversible
Gaussian sources whose real symmetric even trigonometric polynomial spectral
densities have degree at most $m(N-1)$ and are bounded below by
$\gamma I_m$, where $\gamma=\min_c\lambda_{\min}(Q_c)>0$.
Source coordinate variances equal the observed variances. In particular
the maximum delay is $(m-1)(N-1)$, linear in the retained horizon at
fixed channel count.
\end{theorem}
\begin{proof}
It suffices to construct one condition; omit its superscript. Put
$H=N-1$ and $\Delta_{ij}=d_i-d_j=(j-i)H$ for $i<j$.
Define
\[
 S_{ii}(\omega)=v_i+2\sum_{k=1}^{H}C_{ii}(k)\cos(k\omega).
\]
For $i<j$, the numbers $n_k=\Delta_{ij}-k$, $-H\le k\le H$,
are distinct nonnegative integers. Define
\[
 S_{ij}(\omega)=S_{ji}(\omega)
 =\sum_{k=-H}^{H}(2-\mathbf1_{n_k=0})C_{ij}(k)\cos(n_k\omega).
\]
Fourier orthogonality against $d\omega/(2\pi)$ gives source moment
$R_{ij}(k-\Delta_{ij})=C_{ij}(k)$ for each retained $k$.
The coefficient one at frequency zero is required because its integral
is one, whereas $\cos^2(n\omega)$ integrates to $1/2$ for $n>0$.
The diagonal source moments also match the observed auto moments, which
pure delays preserve. Symmetry and stationarity give the corresponding
$j,i$ output entries.

Pointwise, $S_{ii}(\omega)\ge a_i$ and $|S_{ij}(\omega)|\le B_{ij}$. For every real vector $u$, write $|u|$ for its componentwise absolute value. Then
\[
u^{\mathsf T}S(\omega)u
\ge\sum_i a_i u_i^2-\sum_{i\ne j}B_{ij}|u_i||u_j|
=|u|^{\mathsf T}Q|u|
\ge\lambda_{\min}(Q)\|u\|_2^2.
\]
This positive real-even density defines a stationary reversible Gaussian
source. Its degree is at most
$\max_{i<j}(\Delta_{ij}+H)=mH$. Thus its covariances vanish beyond
$mH$, and its Gaussian process is $mH$-dependent. All zero-lag diagonal
moments equal $v_i$. Apply this construction independently to each
condition, preserving the same delays and taking the minimum floor.
\end{proof}
The comparison-matrix condition is unchanged by nonzero rescaling of individual channels. If $D$ is the diagonal scaling matrix, the transformed comparison matrix is $|D|Q_c|D|$, so positive definiteness is equivalent. This condition can hold when a row-dominance inequality fails. For example, white auto-covariances with variances $(1,9)$ and cross lags $(C_{12}(-1),C_{12}(0),C_{12}(1))=(0.02,1,0.01)$ give $Q=\left(\begin{smallmatrix}1&-2.06\\-2.06&9\end{smallmatrix}\right)\succ0$, but its first row is not diagonally dominant. The explicit construction then also proves validity of this retained covariance law. The result concerns this specified sufficient class; the universal theorem also covers targets outside it.


\begin{corollary}[Exact minimum delay span under the comparison-matrix condition]\label{cor:generic-span}
In addition to $Q_c\succ0$ in every condition, suppose that for
 every channel pair $i<j$ and every integer $|\Delta|\le N-2$, some
condition $c$ has
\[
 C^{(c)}_{ij}(\Delta-1)\ne C^{(c)}_{ij}(\Delta+1).
\]
Then the minimum span $\max_i d_i-\min_i d_i$ among pure-delay explanations that use one vector across conditions is exactly $(m-1)(N-1)$. The same minimum is attained by
zero-error strictly positive polynomial sources. For $m=2$, the minimum
absolute relative delay is exactly $N-1$.
The additional inequalities hold outside a finite union of proper
linear subspaces of the retained-moment coordinate space; thus they
hold for Lebesgue-almost-every tuple in this open comparison-matrix class.
\end{corollary}
\begin{proof}
A compatible vector cannot have $|d_i-d_j|\le N-2$, since the specified
condition would violate its retained reflection equation. Hence every
pair of delay values is separated by at least $N-1$. Sort the $m$
delays and add the $m-1$ adjacent gaps to obtain span at least
$(m-1)(N-1)$. The explicit linearly spaced construction attains it.
Each exceptional case fixes a channel pair and relative delay and
requires the same nontrivial linear equality in all conditions. There
are finitely many such cases. Their union has Lebesgue measure zero in
the independent retained coordinates. The positive-definiteness conditions define
an open set, so restricting to it preserves this almost-everywhere
statement. This is a statement about this specified covariance class
and coordinate measure, not a frequency claim about vector autoregressive (VAR) models,
physical processes, or the supplied numerical sample.
\end{proof}
